\begin{recipe}
[% 
    preparationtime = {\unit[10]{min}},
    bakingtime = {\unit[25--30]{min}},
    bakingtemperature = {\protect\bakingtemperature{
        fanoven=\unit[160]{\textcelcius}
    }},
    portion = {\portion{12}},
    %source = {},
    calory = {530kcal | 10g Protein | 52g KH | 30g Fett},
    price = {0,92€},
    created = {12.06.2026},
    %rating = {1},
]
{Kokos-Cashew-Granola}

\graph{% pictures
    big=pic/fruehstueck/kokos-cashew-granola
}

\introduction{%
Knuspriges Granola mit Haferflocken, Cashews, Nüssen, Kokosraspeln und schwarzem Sesam. Süß, nussig, leicht kokoslastig und genau das Zeug, das man eigentlich nur portionieren wollte, bis man dann direkt aus dem Glas weiterisst.
}

\ingredients{
    500 g & Haferflocken\index{Getreide!Haferflocken}\\
    150 g & Cashews\index{Nüsse \& Samen!Cashew}\\
    100--120 g & Nussmischung\index{Nüsse \& Samen!Nüsse}\\
    50--70 g & Kokosraspeln\index{Nüsse \& Samen!Kokos}\\
    25--30 g & Schwarzer Sesam\index{Nüsse \& Samen!Sesam}\\
    \nicefrac{1}{2}--1 TL & Zimt\\
    1 gute Prise & Salz\\
    150 g & Kokosöl\\
    200--210 g & Agavendicksaft
}

\preparation{%
    \step%
    Backofen auf 160 °C Umluft vorheizen. Ein Backblech mit Backpapier auslegen.
    
    \step%
    Haferflocken, Cashews, Nussmischung, Kokosraspeln, schwarzen Sesam, Zimt und Salz in einer großen Schüssel vermengen.
    
    \step%
    Kokosöl vorsichtig schmelzen und mit dem Agavendicksaft verrühren.
    
    \step%
    Die flüssige Mischung zu den trockenen Zutaten geben und alles gründlich vermengen, bis die Haferflocken gleichmäßig überzogen sind.
    
    \step%
    Granola auf dem Backblech verteilen und leicht andrücken.
    
    \step%
    12--15 Minuten backen, dann vorsichtig wenden. Anschließend weitere 10--15 Minuten backen, bis das Granola goldbraun ist.
    
    \step%
    Vollständig auf dem Blech auskühlen lassen. Erst danach auseinanderbrechen und luftdicht lagern.
}

\hint{
    Kokosraspeln bräunen schneller als Haferflocken. Also lieber etwas früher aus dem Ofen nehmen, weil Granola beim Abkühlen noch knuspriger wird.
}

\suggestion[Agavendicksaft]{%
    Wenn du eher 70 g Kokosraspeln nimmst, sind 210 g Agavendicksaft sinnvoller, damit das Granola nicht zu trocken wird.
}

\end{recipe}